\documentclass[10pt]{article}

\usepackage[utf8]{inputenc}
\usepackage[T1]{fontenc}
\usepackage[english]{babel}
\usepackage[left=1.25cm,top=1.05cm,right=1.25cm,bottom=0.9cm]{geometry}
\usepackage{enumitem}
\usepackage{xcolor}
\usepackage{hyperref}
\usepackage{titlesec}
\usepackage{tabularx}
\usepackage{array}
\usepackage{microtype}
\usepackage{lmodern}
\usepackage{needspace}

\setlength{\parindent}{0pt}
\setlength{\parskip}{0pt}
\linespread{1.105}
\pagestyle{empty}

\definecolor{accent}{HTML}{1F4E79}
\definecolor{muted}{HTML}{596675}
\definecolor{rulegray}{HTML}{CAD3DC}

\hypersetup{
    colorlinks=true,
    urlcolor=accent,
    linkcolor=accent,
    pdfauthor={Nguyen Minh Khoa},
    pdftitle={Nguyen Minh Khoa - Full Project CV}
}

\newcommand{\cvsection}[1]{%
    \vspace{6.2pt}%
    {\color{accent}\bfseries\small #1}\par
    \vspace{2.3pt}%
    {\color{rulegray}\hrule height 0.45pt}\vspace{4.0pt}%
}

\newcommand{\entry}[2]{%
    \textbf{#1}\hfill {\color{muted}\small #2}\par
}

\newcommand{\subentry}[1]{%
    {\color{muted}\small #1}\par
}

\setlist[itemize]{
    leftmargin=13pt,
    labelsep=5pt,
    itemsep=2.8pt,
    topsep=2.8pt,
    parsep=0pt,
    partopsep=0pt
}

\begin{document}

\begin{center}
    {\LARGE\textbf{NGUYEN MINH KHOA}}\\[-1pt]
    {\color{accent}\textbf{AI Engineer \& Researcher}}\\[3pt]
    {\small Ho Chi Minh City, Vietnam \textbullet\
    \href{tel:+84907713819}{+84 907 713 819} \textbullet\
    \href{mailto:ngminhkhoayj2706@gmail.com}{ngminhkhoayj2706@gmail.com} \textbullet\
    \href{https://www.linkedin.com/in/ngminhkhoa1410}{LinkedIn} \textbullet\
    \href{https://github.com/Purin1410}{GitHub} \textbullet\
    \href{https://scholar.google.com/citations?user=tCq7yoQAAAAJ}{Google Scholar} \textbullet\
    \href{https://purin1410.github.io}{Portfolio}}
\end{center}

\vspace{-2pt}

\cvsection{PROFILE}
Applied AI engineer and Research Assistant at AiTA Lab, FPT University, working across model experimentation, data pipelines, and inference systems. Research spans handwritten mathematical expression recognition and text-to-molecule generation, with a current focus on AI for intelligent manufacturing.

\cvsection{RESEARCH EXPERIENCE}
\entry{Research Assistant | AiTA Lab, FPT University}{Oct 2024 - Present}
\begin{itemize}
    \item Developed and evaluated deep-learning methods across HMER and text-to-molecule generation, covering data preparation, model training, ablation studies, and reproducible experiment tracking.
    \item Selected published outcomes include up to \textbf{+5 ExpRate points} for Mask CoMER, \textbf{+4 ExpRate points} for MTL + CBARM, and \textbf{+3.4 BLEU} with approximately \textbf{30\% lower Levenshtein distance} for ChemAligner-T5.
    \item \textbf{5 papers across three research topics:} HMER (\textit{Mask CoMER}; \textit{MTL + CBARM}), text-to-molecule generation (\textit{LexiChem}; \textit{ChemAligner-T5}), and active learning with vision-language models (\textit{CorrTie}).
\end{itemize}

\cvsection{PROJECTS}
\Needspace{16\baselineskip}
\entry{LexiChem | AI \& Inference Contributor}{2025 - 2026}
\subentry{Capstone combining text-to-molecule research with a complete local demo for molecule generation and exploration.}
\begin{itemize}
    \item Worked in a 3-person capstone team; owned the AI data pipeline, SMILES/SELFIES conversion, method development, ablations, and the inference path between research models and the application backend.
    \item Combined L+M-24, Mol-Instructions, and ChEBI-20; used RDKit to parse structures, validate valence, canonicalize SMILES, remove duplicate molecules, and convert cleaned structures to SELFIES for training.
    \item Developed the shared-latent alignment method and ran controlled ablations. On ChEBI-20, improved BioT5+ from \textbf{87.2 to 88.8 BLEU}, \textbf{52.2 to 52.9 exact match}, and \textbf{90.7 to 93.5 MACCS FTS}, while maintaining \textbf{100\% validity}.
    \item Wrote the paper Introduction and full Results \& Discussion, and co-wrote Methods; presented LexiChem at APWeb-WAIM 2026 in Da Nang.
    \item Integrated the Experiments workflow with FastAPI and NVIDIA Triton, including prompt normalization, model dispatch, RDKit validation of returned SMILES, and experiment persistence in MongoDB.
    \item Worked with the team to productize generation into a local application with Dashboard, Experiments, Simulation, 2D/3D molecular exploration, and docking-oriented workflows. RAG/Knowledge remains in development.
\end{itemize}

\Needspace{14\baselineskip}
\entry{NexOps | AI, Simulation \& Platform Lead}{2025 - Present}
\subentry{Local-first MES/SCADA prototype linking factory simulation, production workflows, dashboards, and real-device control.}
\begin{itemize}
    \item Led the AI, simulation, and platform work for a working local prototype covering machine telemetry, production planning, work orders, alarms, event history, and OEE by shift.
    \item Built repeatable machine and shift scenarios so planning, alarms, downtime, and OEE can be exercised through the same deterministic workflow during development and demos.
    \item Designed the telemetry path from simulator and edge events through MQTT/EMQX into FastAPI services, TimescaleDB, and operator-facing dashboard state.
    \item Connected the local runtime to production-facing views for machine status, planning, alarm handling, and OEE reporting rather than leaving the simulator as a standalone data generator.
    \item Integrated a Bambu Lab printer over LAN and exercised real-device control as a separate hardware path inside the prototype.
\end{itemize}

\Needspace{12\baselineskip}
\entry{HMER Research | Research Contributor}{2024 - 2026}
\subentry{Research track on converting handwritten mathematical expressions into structured LaTeX, with published work on pretraining and multi-task attention refinement.}
\begin{itemize}
    \item Contributed to \textit{Mask CoMER}, a two-stage HMER approach using masked-language pretraining followed by task training with stochastic-depth regularization; published results improved ExpRate by up to \textbf{5 points} over the CoMER baseline.
    \item Contributed multi-task learning experiments for the \textit{MTL + CBARM} line, where the main LaTeX-sequence task is trained with a structural auxiliary task and attention refinement; published results improved ExpRate by approximately \textbf{4 points} over the CoMER baseline.
    \item Trained and evaluated PyTorch/PyTorch Lightning models on CROHME benchmarks and tracked experimental runs with Weights \& Biases.
    \item Worked across model experimentation, evaluation, and reproducibility rather than only manuscript preparation; both research lines resulted in published papers.
\end{itemize}

\Needspace{14\baselineskip}
\entry{TrendRadar | Data Pipeline Owner}{2025}
\subentry{AI-driven trend analytics platform for short-form content, built to turn heterogeneous TikTok, TikTok Shop, and Reels data into normalized records for downstream analytics and ML.}
\begin{itemize}
    \item Designed and implemented the end-to-end ingestion and preprocessing pipeline from data collection through normalization and PostgreSQL storage.
    \item Built asynchronous ingestion with Pydantic schema validation to enforce consistent contracts across heterogeneous sources.
    \item Added idempotent processing, deduplication, retry/backoff, and quarantine handling for malformed records to stabilize scheduled runs.
    \item Added pytest coverage, structured logging/monitoring hooks, and Docker-based local execution for reproducible pipeline development.
    \item Reduced ingestion failures by \textbf{30\%} and authored \textbf{8 typed schemas} for source validation and normalization.
\end{itemize}

\Needspace{14\baselineskip}
\entry{Research Ops | Developer}{2026 - Present}
\subentry{Personal pre-alpha workspace for keeping sources, evidence, experiments, revisions, and human decisions traceable across long-running research cycles.}
\begin{itemize}
    \item Designed the project around reducing human attention required to keep trustworthy research moving while preserving provenance, checkpoints, and researcher judgment.
    \item Implemented a local research library, PDF/citation handling, work sessions, evidence links, manuscript exports, and resumable study state backed by FastAPI and PostgreSQL.
    \item Linked source files and cited spans through SHA-256 content digests so evidence remains tied to the exact source version used for a decision or manuscript revision.
    \item Added revision-aware exports and approval gates for runtime changes; the researcher retains control of direction, interpretation, and scientific conclusions while agents handle bounded execution.
    \item Current status is explicitly pre-alpha: no public deployment or measured long-term benefit yet, and model execution/GPU evaluation remain unfinished.
\end{itemize}

\cvsection{PUBLICATIONS}
\textbf{LexiChem: Shared-Latent Alignment for Faithful Text-to-Molecule Generation in SELFIES Space}\\[-1pt]
{\small APWeb-WAIM 2026 | Accepted and presented | \textbf{Scopus-indexed proceedings}}\par
\vspace{2.8pt}
\textbf{CorrTie: Correction-aware Tie-breaking for Active Learning with Vision-Language Models}\\[-1pt]
{\small APWeb-WAIM 2026 | Accepted | \textbf{Scopus-indexed proceedings}}\par
\vspace{2.8pt}
\textbf{Enhancing handwritten mathematical expression recognition with convolutional block attention refinement module and multi-task learning}\\[-1pt]
{\small Journal of Information and Telecommunication | \textbf{SCImago Q2 (2025)} | Published online Aug 2026}\par
\vspace{2.8pt}
\textbf{ChemAligner-T5: A Unified Text-to-Molecule Model via Representation Alignment}\\[-1pt]
{\small ICCIES 2026 | Springer CCIS proceedings | Published}\par
\vspace{2.8pt}
\textbf{Mask CoMER: Enhancing Handwritten Mathematical Expression Recognition with Masked Language Pretraining and Regularization}\\[-1pt]
{\small ICDAR 2025 | \textbf{ICORE 2026 Rank A} | Published}

\Needspace{8\baselineskip}
\cvsection{EDUCATION}
\entry{Bachelor of Science in Artificial Intelligence | FPT University, HCMC}{2022 - 2026}
{\small Academic requirements completed; degree conferral expected Nov 2026.}

\cvsection{AWARDS \& SCHOLARSHIPS}
\begin{itemize}
    \item \textbf{50\% Tuition Scholarship}, FPT University (2022), awarded through the university scholarship examination and covering 50\% of tuition for the full programme.
    \item \textbf{RESFES HCMC Second Prize}, ChemAligner-T5 (2026).
    \item \textbf{RESFES HCMC Second Prize}, Mask CoMER (2025).
    \item \textbf{Demo Pitching Day}, VND 50M team award/project funding (2025).
    \item \textbf{Innovation Quest}, VND 40M team award/funding for TrendRadar / FrontAI (2025).
    \item \textbf{HCMC Open April Olympiad}, Silver Medal in Mathematics (2021).
\end{itemize}

\cvsection{TECHNICAL SKILLS}
{\small
\textbf{AI \& Research:} Python, PyTorch, PyTorch Lightning, Hugging Face Transformers, scikit-learn, OpenCV, RDKit, Weights \& Biases, MLflow, self-supervised learning, vision-language models, cross-modal learning\\
\textbf{Backend \& Data:} FastAPI, Pydantic, Pandas, PostgreSQL/TimescaleDB, MongoDB, Redis, MQTT/EMQX, REST, ETL/data pipeline design\\
\textbf{Engineering \& Reproducibility:} Docker, Docker Compose, Linux, Git, Nginx, Playwright, LaTeX, scientific writing
}

\cvsection{CERTIFICATIONS}
\begin{itemize}
    \item \textbf{AIO 2024 | AI VIET NAM}: Machine Learning; Deep Learning; Computer Vision \& NLP; GenAI \& LLMs.
    \item \textbf{Machine Learning Specialization}, Stanford Online / DeepLearning.AI (2023).
    \item \textbf{Deep Learning Specialization}, DeepLearning.AI (2024).
    \item \textbf{Natural Language Processing Specialization}, DeepLearning.AI (2025).
    \item \textbf{DeepLearning.AI TensorFlow Developer Professional Certificate} (2025).
    \item \textbf{Big Data Specialization}, University of California San Diego (2024).
    \item \textbf{AI Foundations for Everyone Specialization}, IBM (2023).
    \item \textbf{Introduction to AI}, Google (2026).
    \item \textbf{IBM Full Stack Software Developer Professional Certificate} (2024).
    \item \textbf{AWS Cloud Technical Essentials}, Amazon Web Services (2025).
    \item \textbf{Software Development Lifecycle Specialization}, University of Minnesota (2024).
    \item \textbf{Research Methodologies}, Queen Mary University of London (2025).
    \item \textbf{Being a Researcher in Information Science and Technology}, Politecnico di Milano (2025).
    \item \textbf{Understanding Research Methods}, University of London / SOAS University of London (2025).
    \item \textbf{Introduction to Research for Essay Writing}, University of California, Irvine (2025).
    \item \textbf{Advanced Writing}, University of California, Irvine (2025).
    \item \textbf{Academic Skills for University Success Specialization}, The University of Sydney (2023).
    \item \textbf{Project Management Principles and Practices Specialization}, University of California, Irvine (2025).
    \item \textbf{Ethics in the Age of AI Specialization}, LearnQuest (2024).
\end{itemize}

\end{document}
